% SPDX-License-Identifier: Apache-2.0
\section{A Pipeline Is an Async Function}
\label{sec:e-pipeline}

A pipeline and a sequence are one construct driven two ways.

A sequence runs one invocation to completion before starting the next. A
pipeline starts a new invocation every cycle, so several are in flight at
once, each at a different point in the body. Nothing about the body
differs. The difference is the rate at which invocations start, and that
belongs at the call site.

\begin{lstlisting}[caption={A pipeline, a sequence and a parallel process, all at once. Nobody placed the stage boundary: it is where an operator with latency is awaited.},label={lst:pipeline}]
use crate::pipeline::{add, mul};

async fn mac(a: U<32>, b: U<32>, prev: U<64>) -> U<64> {
    let p = mul(a, b).await;   // the multiplier's latency
    add(prev, p).await         // and the adder's
}
\end{lstlisting}

\subsection{Nobody places the stage boundary}

An earlier draft of this listing wrote \code{tick().await} to say where
the register goes. That was wrong, and wrong in a way the project has
already ruled on: the eighth thesis~\cite{theses} asks the language to
stop making a designer place registers by hand, and a hand-placed tick is
exactly that under a new name.

Operators that have latency are \code{async} instead, and they live in
\code{crate::pipeline}, because each of them is a pipeline stage the
moment it is awaited. Combinational operators live in
\code{crate::funcs}, so the module a design imports from already says
whether the operation can cost time. A multiply takes as long as the
mapping decides, so
\code{mul} is awaited and the caller never learns the number. An adder is
the same: \code{add} is awaited too, because whether an addition costs a
cycle is a question about the adder and the clock, not about the design.

\subsection{Why there is no \code{impl Add}}

Writing \code{prev + p} would be shorter, and it is not offered.

An operator has to return a value. An operation whose latency the mapping
decides cannot, so \code{+} could only be provided for the zero-latency
case, and a design that wrote it would have chosen a combinational adder
without saying so. The shorter spelling would be the one that quietly
makes the decision the language exists to defer.

The objection to this is verbosity, and it rests on a misreading.
\code{.await} is not a claim that a cycle is spent. It is a refusal to
decide here. An adder the mapping gives zero latency is still awaited:
the netlist makes it a wire, and the design that awaited it did not
pay for the question. The run does not follow the mapping, and its
\code{add} spends one cycle.

The pipeline's depth then follows from the operators the design used. Two
multiplies in sequence give two boundaries, and rebinding \code{mul} to a
single-cycle implementation removes them, without an edit to
Listing~\ref{lst:pipeline}.

Three constructs go, and one of them takes a compiler pass with it.

\subsection{The pipe declaration disappears}

LHDL's \code{pipe} declared a value that crosses a stage boundary, so that
live-range analysis would know to hold it. TxHDL had the same notion under
a different name.

In an async function, a local that is live across an \code{.await} is held
in the state machine the compiler generates, because that is what an async
function is. \code{p} in Listing~\ref{lst:pipeline} crosses the await inside
\code{mul} and needs no declaration. The analysis the \code{pipe} keyword existed to feed
is one that \code{rustc} performs already.

\subsection{What this does not settle}

Three things, and none of them is free.

\begin{itemize}
  \item \term{Named stages are gone.} A boundary is an await, not a named
        block. A name becomes a comment.
  \item \term{The initiation interval moves to the call site.} It was
        implied by the choice of keyword and is now an argument where the
        body is driven. That is a fair trade and not a saving.
  \item \term{Resource sharing has to be proved.} A real pipeline shares
        one multiplier across every item in flight, while several
        concurrent invocations each appear to want their own. They are the
        same code at the same await index, so the lowering can share them.
        That is a claim about the compiler being written, not about Rust,
        and it is the largest piece of unfinished work in this proposal.
\end{itemize}

One thing improves. TxHDL had a \code{stall} statement for a stage that
cannot proceed. A stall is now an await that has not completed, which is
the mechanism rather than a construct layered on top of it.

\code{mac} does not loop, and should not. It is an operation, invoked
once per item; the unit that calls it is what loops, per
Section~\ref{sec:e-regread}, and awaiting \code{mac} inside that loop is
what makes each call one item in flight.

\subsection{It has been run}

The claim above was written before anything ran it. The runtime now
has \code{pipeline::drive}, which starts an invocation of an
\code{async fn} at every edge at which an input is offered and polls
every invocation in flight as a process of its own, and the examples
document runs a two-await \code{mac} through it. At every edge from
the second, one invocation is in its multiplier and another in its
adder; from the third, a result lands every cycle. The stage boundaries
were the awaits, and nobody placed them. Resource sharing, the item
above that had to be proved, is still not: the prototype gives each
invocation its own multiplier, and it is the lowering that would share
one.
